\section{Sobre el Tema}

\subsection{Palabras Clave}

\begin{itemize}
  \item \textbf{Reclutamiento Digital:} Modalidad relativamente moderna para la manipulación y vinculación de menores a grupos o redes de crimen organizado desde entornos digitales tales como: redes sociales, videojuegos u otras plataformas.
  \item \textbf{Ciudadanía Digital:} Según la UNESCO, este término comprende habilidades que sobrepasan el conocimiento del uso de dispositivos tecnológicos. Comprende el uso que se les da y que este sea ético y crítico, el ejercicio de los derechos en plataformas digitales tales como todos aquellos en relación a la información y los datos, la alfabetización mediática e informacional y la convivencia y el bienestar. UNICEF considera la huella digital como un término en evolución constante y que incorpora un conjunto de competencias que habilitan a los ciudadanos a navegar la red y hacer un uso correcto de la misma, bajo el conocimiento de las posibilidades así como los riesgos y las soluciones a dichos retos.
  \item \textbf{Huella Digital:} Este concepto se refiere al rastro de la colección de datos que deja el uso de internet ya sea de uso personal o corporativo. Estos datos pueden ser dados con consentimiento directo o indirecto y el espectro de rastreo de actividad es inmenso incluyendo tiempo de uso, contenido consumido e historial de navegación e incluso datos personales.
  \item \textbf{Cookies:} Este es el término común utilizado para referirse a los archivos que se almacenan en el dispositivo del usuario de internet al visitar una página web. Bajo esta definición, cada operador del sitio web determina qué información se obtiene a través de esta herramienta y cuál es su propósito. Según la entidad gestora pueden existir distintos tipos como por ejemplo:
  \begin{itemize}
    \item Cookies propias
    \item Cookies de terceros
    \item Cookies de sesión
    \item Cookies persistentes
    \item Cookies tecnicas
    \item Cookies de preferencia
    \item Cookies de análisis
    \item Cookies de publicidad comportamental
  \end{itemize}
\end{itemize}

Estas tienen validez de carácter similar al contractual y, por ende, dificultan o incluso imposibilitan la apelación en contra del almacenamiento de los datos estipulados en dicho archivo presentando así un peligro que podría llegar a considerarse legalmente protegido.

\begin{itemize}
  \item \textbf{Extremismo Violento:} Según la definición de la ONU, el extremismo violento no tiene más que una definición ambigua, dada su multidimensionalidad. Suele basarse en la intolerancia a un cierto aspecto social, cultural, etnico o religioso, cuyo método de dispersión principal son los medios digitales públicos.
  \item \textbf{Ciberacoso:} Nueva modalidad de acoso o intimidación llevada a cabo a través de redes sociales, juegos o incluso plataformas de comunicación en dispositivos tecnológicos. Este se basa en un comportamiento repetitivo con el objetivo de asustar, enfurecer o humillar a la víctima.
\end{itemize}

Suele ser uno de los componentes del acoso, pero el ciberacoso deja un registro digital fundamental para la detección del abuso especialmente en casos de menores.

\begin{itemize}
  \item \textbf{Grooming:} El grooming es el término utilizado para referirse a formas delictivas de acoso que implican el contacto de una persona adulta con un niño, niña o adulescente para adquirir un cierto nivel de confianza con el menor con el objetivo de aislar al menor de su red de apoyo generando secretismo e intimidad. De esta manera, el adulto se aprovecha del vínculo creado con el menor con el objetivo de involucrar a este por lo general en algún tipo de actividad sexual o, más adecuadamente para el tema tratado en redes de captación de menores a través de la coacción.
  \item \textbf{Comunidades Peligrosas:} Grupos o comunidades en línea de caracter social, educativo o de ocio con objetivo de fomentar la radicalización a través de la promoción de la intolerancia, discursos de odio, la discriminación y la violencia por distintos aspectos ya sean religiosos, etnicos, etc. También se puede asociar con la accesibilidad a problemas de salud o incluso el abuso sexual de menores.
  \item \textbf{Abuso y Explotación Sexual Infantil en Línea (AESI):} Este es el término general que engloba todos los actos de violencia sexual hacia los menores a través de plataformas digitales. El ICMEC y Save the Children incluyen aquí todas las formas de explotación que no tiene por qué ser agresión y explotación física o psicológica únicamente. Se incluyen todas las formas de beneficio del delito, todo tipo de manipulación o engaño, extorsión etc. 3
  \item \textbf{Sharenting:} Según Stacey Steinberg, el Sharenting es la publicación de los hijos en redes o plataformas digitales así como información sobre ellos, exponiéndoles a todos los peligros que eso puede llegar a conllevar.
  \item \textbf{Sextorsión / sextorsión ffinanciera:} Una forma de extorsión en la que una persona amenaza con divulgar imágenes, videos o mensajes de contenido íntimo o sexual de otra persona a cambio de dinero, más imágenes, favores sexuales o cualquier otro beneficio. En el caso de sextorsión financiera, el objetivo principal es obtener dinero.
  \item \textbf{Deep Fakes:} Contenido audiovisual generado o manipulado mediante inteligencia artificial para crear representaciones falsas pero altamente realistas de personas, simulando su apariencia, voz o acciones. Pueden utilizarse con fines legítimos, pero también para difundir desinformación, cometer fraudes, dañar la reputación de una persona, entre otros.
  \item \textbf{Hashing:} Proceso criptográfico que convierte datos en un valor de longitud fija (hash) para verificar la integridad de la información sin revelar su contenido original.
  \item \textbf{End-to-end encryption:} Un mecanismo de seguridad que cifra la información desde el dispositivo del emisor hasta el del destinatario, asegurando que únicamente las partes autorizadas puedan acceder al contenido de la comunicación.
\end{itemize}

\subsection{Introducción al tema}

En un mundo de incertidumbre tal y como en el que nos encontramos hoy en día dada la evolución constante de las tecnologías, los derechos de todos los niños y niñas del mundo se encuentran en una posición de alta vulnerabilidad. La radicalización y la explotación digital se nutren de un entorno digital repleto de datos cada vez a mayor disposición que aún no se ha coordinado con la educación necesaria para obtener los resultados positivos que prometen. A pesar de ser herramientas de valor incalculable en cuanto a comunicación, desarrollo y trabajo, el aumento de datos en la nube suponen un acceso de dimensiones alarmantes a la vida y las mentes de los niños, niñas y adolescentes.

El grooming, a pesar de su definición tradicional focalizada en el abuso sexual de los menores mediante plataformas dígitales, a día de hoy toma un significado muchisimo mas amplio, enfocado en areas ideológicas, extremistas y criminales mediante la coacción, el miedo y el secretismo. A través de este tipo de prácticas, se dan dos fenómenos complementarios trascendentales siendo estos la radicalización y la captación.

A pesar de ser un facilitador, la tecnología no es el origen, sino un magnificador que ha facilitado el alcance de los daños que pueden causar este tipo de delitos de abuso y explotación, como explicó Gordon Alexander, exdirector de la Oficina de Investigación de la UNICEF.\footnote{UNICEF. (s. f.). El Estado Mundial de la Infancia 2017: Los niños en un mundo digital — Resumen. \url{https://www.unicef.org/media/48611/file/SOWC_2017_Summary_SP.pdf}} Es por esto que es imperativo comprender las diferentes facetas, características y orígenes del problema real a la hora de debatir sobre métodos para mitigar y eliminar las causas y los efectos que la violencia, la radicalización y la explotación de menores presentan.La brecha jurídica en el marco legislativo internacional es el mayor reto que este comité ha de afrontar combinando la cooperación y la coordinación educativa y legislativa respetando la privacidad de los datos y la decisión individual.

El grooming, la adoctrinación y la influencia sobre los menores es cada vez mayor y constituye un reto multidimensional y transnacional. La violencia, la explotación y la radicalización de menores requiere de una respuesta eficaz, híbrida, multidimensional y multilateral para poder garantizar la protección de los derechos de los niños, niñas y adolescentes de todo el mundo.

Dada la presencia temprana del teléfono móvil en los niños, según UNICEF la media de edad para el acceso al móvil es a los 10,8 años, donde el 41\% de los niños y niñas disponen de móvil propio y continúa subiendo según aumenta el rango de edad.\footnote{UNICEF España. (s. f.). Nuevo informe sobre el impacto de la tecnología en la infancia y la adolescencia. \url{https://www.unicef.es/noticia/nuevo-informe-sobre-el-impacto-de-la-tecnologia-en-la-infancia-y-la-adolesc} encia} El acceso tan temprano a los dispositivos móviles y las plataformas digitales, en combinación con la falta de adaptación en los programas educativos, presentan un riesgo en incremento. Aproximadamente un 12.6\% global de los menores ha sufrido de publicación y difusión de imágenes de connotación sexual sin consentimiento y casi uno de cada diez niños pueden esperar ser víctimas de sextorsión.\footnote{Universidad de Edimburgo. (s. f.). Un informe muestra la magnitud del abuso infantil en Internet. \url{https://www.ed.ac.uk/news/report-shows-scale-of-online-abuse-of-children}}

En cuanto a la radicalización de menores a mediante plataformas digitales, el proyecto de Disrupting Harm que proyecta un esfuerzo conjunto de UNICEF, INTERPOL y EPCAT estipula que hasta un 19\% de menores de entre 12 y 17 años sufrieron algún tipo de explotación digital en 12 meses, utilizando como plataformas de captación aplicaciones de mensajería y un 12\% a través de videojuegos en línea.\footnote{UNICEF Innocenti. (s. f.). Disrupción del daño. \url{https://www.unicef.org/innocenti/projects/disrupting-harm}} \footnote{UNICEF Montenegro. (s. f.). Informe sobre derechos digitales en Montenegro. \url{https://www.unicef.org/montenegro/en/media/26231/file/DH-Montenegro-Report-eng-FINAL.pdf.pdf}} Por último, el UN CTED concluyó que las estrategias de reclutamiento atacan principalmente entre menores de 8 a 9 años de edad y el Soufan Center en 2026 calculó que los menores de edad conforman el 20\% de los arrestos por delitos en relación al terrorismo en jurisdicciones de Europa y aliados.\footnote{The Soufan Center. (s. f.). Lucha contra la radicalización de los jóvenes. \url{https://thesoufancenter.org/research/combating-youth-radicalization/}}

\subsection{Contexto e historia}

Tras los inicios de la era tecnológica en 1940, la creación de internet durante la década de los 1990 y la masificación de los smartphones y las redes sociales a partir de los años 2000, el acceso a internet se convirtió en un aspecto crucial para cualquier nación desarrollada, encontrando su sitio en cada una de nuestras vidas, incluyendo las vidas de nuestros niños y niños.\footnote{Instituto de Ingeniería y Tecnología. (s. f.). ¿Cómo empezó la revolución digital en el mundo? \url{https://utec.edu.pe/blog-de-carreras/administracion-y-negocios-digitales/como-empezo-la-revolucion-digit} al-en-el-mundo} Así mismo, el crimen ha encontrado la forma de impregnarse en estas plataformas ya sea en cuanto a reclutamiento de menores alentando la radicalización y la pérdida de pensamiento crítico o como acceso directo a la sexualización y la objetivización de los menores. Más tarde el rápido desarrollo de la inteligencia artificial se ha coordinado de forma alarmante con el crimen, mientras que la educación y el avance legislativo se encuentran en una parálisis de desarrollo debido principalmente al Big Data.\footnote{World Economic Forum. (2025, enero). Avances tecnológicos y desarrollo humano: Una historia de dos mundos. \url{https://es.weforum.org/stories/2025/01/avances-tecnologicos-y-desarrollo-humano-una-historia-de-dos-m} undos/}

El marco internacional está en un punto de descoordinación que de ninguna forma favorece la protección de los derechos de los niños, presentando tres bloques principales de actuación. En primer lugar encontramos la estrategia centrada en la regulación tecnológica y por ende restricción parcial de acceso a los datos y al uso que se les da a estos que, por ende, restringe el desarrollo tecnológico ininterrumpido. La Unión Europea es el mejor ejemplo de este tipo de estrategia en documentos como la Ley de Inteligencia Artificial de la UE de 2024.\footnote{Comisión Europea. (2024). Reglamento (UE) 2024/1689 del Parlamento Europeo y del Consejo, de 13 de junio de 2024, por el que se establecen normas armonizadas en materia de inteligencia artificial (Reglamento de Inteligencia Artificial). EUR-Lex. \url{https://eur-lex.europa.eu/legal-content/EN/TXT/?uri=CELEX\%3A32024R1689}}

En segundo lugar, existe otro método de regulación centrado principalmente en los usos de la tecnología. Y por último la regulación de la infraestructura tecnológica.

Evidentemente, estos marcos no garantizan el salvaguardo de los derechos de los niños, y en la mayoría de los casos, ni siquiera se tienen en cuenta. Ese es vuestro objetivo como representantes de los estados miembros, asegurar una coordinación transnacional que mediante la cooperación, el debate y el respeto se asegure de mitigar todos los riesgos y las consecuencias ya existentes y por existir, o al menos intentarlo, de manera que la correlación entre el desarrollo y el crimen sea revertida e idealmente, erradicar el radicalismo y la explotación a través de los dispositivos tecnológicos y todas sus consecuencias en la infancia.
